\documentclass[12pt]{article}
\usepackage{amsmath,amsthm,amssymb}
\usepackage{hyperref}
\usepackage{geometry}
\usepackage{booktabs}
\geometry{margin=1in}

\newtheorem{theorem}{Theorem}[section]
\newtheorem{proposition}[theorem]{Proposition}
\newtheorem{lemma}[theorem]{Lemma}
\newtheorem{corollary}[theorem]{Corollary}
\newtheorem{definition}[theorem]{Definition}
\newtheorem{remark}[theorem]{Remark}

\DeclareMathOperator{\Spec}{Spec}
\DeclareMathOperator{\Frob}{Frob}
\newcommand{\ZZ}{\mathbb{Z}}
\newcommand{\RR}{\mathbb{R}}
\newcommand{\CC}{\mathbb{C}}
\newcommand{\FF}{\mathbb{F}}
\newcommand{\QQ}{\mathbb{Q}}
\newcommand{\Ql}{\mathbb{Q}_\ell}
\newcommand{\Hminus}{\mathcal{H}_-}
\newcommand{\MRH}{M_{\mathrm{RH}}}
\newcommand{\PAC}{P_{\mathrm{AC}}}

\title{An Unsatisfiability Certificate for Geometric and Spectral\\
Proof Strategies for the Riemann Hypothesis}

\author{J.~Xie\thanks{Independent researcher. Contact: anatomy1602@gmail.com}}
\date{May 2026}

\begin{document}
\maketitle

% ---- ABSTRACT ----
\begin{abstract}
We apply the meta-methodological framework of unsatisfiability certificates
to the problem of proving the Riemann Hypothesis (RH) via geometric or
spectral methods. The framework, developed in prior work, identifies three
conditions under which a class of proof strategies is provably closed: a
\emph{discriminating property} $P$ that (1) any successful proof must
satisfy, (2) provably prevents success, and (3) is not decidable within
the proof class itself (self-referential safety).

We construct an explicit unsatisfiability certificate $(M_{\mathrm{RH}},
f_{\mathrm{RH}}, P_{\mathrm{AC}})$ where $M_{\mathrm{RH}}$ is the class
of geometric and spectral proof strategies for RH, and $P_{\mathrm{AC}}$
(arithmetic compactness) is the discriminating property. The conflict
condition follows from Deligne's proof of the Weil conjectures (1974):
any compact arithmetic scheme $X$ whose $L$-function equals the Riemann
zeta function would force all zeta zeros onto the critical line---which
is RH itself. The self-referential safety condition follows from the fact
that deciding $P_{\mathrm{AC}}$ requires establishing this identification,
which is outside the scope of any strategy in $M_{\mathrm{RH}}$.

The certificate implies: any proof of RH must use a strategy outside
$M_{\mathrm{RH}}$. This is not a proof that RH is unprovable---by
Shoenfield absoluteness, RH is not independent of ZFC. It is a precise
structural boundary on two of the most studied proof paradigms.
\end{abstract}

\tableofcontents
\newpage

% ---- SECTION 1 ----
\section{Introduction}
The Riemann Hypothesis (RH) asserts that all non-trivial zeros of
$\zeta(s)$ lie on the critical line $\mathrm{Re}(s) = \tfrac{1}{2}$.
Two broad proof paradigms have attracted sustained attention.

\textbf{Geometric strategies.} Weil proved RH for curves over finite
fields using the Hodge Index Theorem on $C \times C$. The natural
question is whether an arithmetic analogue exists for $\Spec(\ZZ)$.
Arakelov geometry (Faltings 1984), $\mathbb{F}_1$ geometry
(Connes--Consani 2014), and tropical geometry have all been proposed
as frameworks.

\textbf{Spectral strategies.} The Hilbert--P\'olya conjecture proposes
that the imaginary parts $\gamma_n$ of the zeta zeros are eigenvalues
of a self-adjoint operator. Connes (1999) gave a spectral interpretation
via the scaling operator on $L^2(C_\QQ)$.

Prior work (Phases 17--19 of this research program) identified specific
obstacles within each known framework. However, as correctly noted by
reviewers, listing obstacles in known frameworks does not constitute a
proof that \emph{all} geometric or spectral strategies fail---new
frameworks might avoid the listed obstacles.

This paper closes that gap by constructing an \emph{unsatisfiability
certificate} in the sense of the meta-methodological framework
\cite{metamethod2026}. The certificate covers the entire class
$\MRH$ of geometric and spectral strategies, not just known instances.
The key tool is Deligne's theorem \cite{deligne1974}, which forces
any compact arithmetic scheme to satisfy RH for its own $L$-function---
making the identification with $\zeta(s)$ circular.

\subsection*{Main result}

\begin{theorem*}[Informal]
The triple $(\MRH, f_{\mathrm{RH}}, \PAC)$ is an unsatisfiability
certificate: no geometric or spectral proof strategy (formalized in
ZFC, not assuming RH) can prove RH.
\end{theorem*}

\subsection*{What this result does not say}

By Shoenfield absoluteness \cite{shoenfield1961}, RH is a
$\Pi^0_1$ statement and is therefore not independent of ZFC: it is
either provable or disprovable. The main result does not contradict
this---it asserts only that the proof, if it exists, cannot proceed
via geometric or spectral methods as defined in Section~\ref{sec:model}.

\subsection*{Organization}

Section~\ref{sec:framework} recalls the meta-methodological framework.
Section~\ref{sec:weil} states the Weil positivity criterion.
Section~\ref{sec:model} defines $\MRH$ precisely.
Section~\ref{sec:pac} defines $\PAC$ and verifies the three conditions.
Section~\ref{sec:cert} states and proves the main theorem.
Section~\ref{sec:second} applies the Second Law.
Section~\ref{sec:discussion} discusses scope and open questions.

% ---- SECTION 2 ----
\section{The Meta-Methodological Framework}
\label{sec:framework}

We recall the framework from \cite{metamethod2026}. The framework
identifies a structural condition---\emph{self-referential safety}---
that governs the success and failure of impossibility proofs across
multiple domains (circuit complexity, mathematical logic, and, in this
paper, proof strategies for RH).

\begin{definition}[Proof strategy model]
A \emph{proof strategy model} is a pair $M = (S, C)$ where $S$ is a
class of candidate proof strategies and $C$ is a set of constraints
every strategy must satisfy.
\end{definition}

\begin{definition}[Discriminating property]
\label{def:discrim}
Given a model $M = (S,C)$ and a target $f: S \to \{0,1\}$ with ideal
value $v^* = 1$, a \emph{discriminating property} is a predicate $P$
on strategies such that:
\begin{enumerate}
\item[\textup{(Necessity)}] Any $\mathcal{A} \in M$ with
  $f(\mathcal{A}) = 1$ must satisfy $P(\mathcal{A})$.
\item[\textup{(Conflict)}] $P(\mathcal{A})$ provably implies
  $f(\mathcal{A}) = 0$.
\end{enumerate}
\end{definition}

\begin{definition}[Self-referential safety]
\label{def:safety}
A discriminating property $P$ is \emph{self-referentially safe} with
respect to $M$ if no strategy $\mathcal{A} \in M$ can decide $P$:
\[
\nexists\, \mathcal{A} \in M \text{ such that } \mathcal{A}
\text{ decides } P.
\]
\end{definition}

\begin{definition}[Unsatisfiability certificate]
A triple $(M, f, P)$ is an \emph{unsatisfiability certificate} if $P$
is a discriminating property that is self-referentially safe with
respect to $M$.
\end{definition}

\begin{theorem}[{\cite[Theorem~2.6]{metamethod2026}}]
\label{thm:framework}
If $(M, f, P)$ is an unsatisfiability certificate, then no strategy
in $M$ achieves $f = 1$.
\end{theorem}

\begin{proof}
By Necessity, any strategy achieving $f=1$ satisfies $P$.
By Conflict, any strategy satisfying $P$ has $f=0$. Contradiction.
\end{proof}

\begin{remark}
Self-referential safety does not appear in the logical proof above.
Its role is \emph{methodological}: it explains why the discriminating
property $P$ can be identified and used without the proof class
being able to simulate or circumvent it. In the circuit complexity
setting, this corresponds to the Natural Proofs barrier
\cite{razborov-rudich1994}. In the present setting, it corresponds
to the fact that deciding $\PAC$ requires establishing an identification
that is equivalent to RH itself (Section~\ref{sec:pac}).
\end{remark}

% ---- SECTION 3 ----
\section{The Weil Positivity Criterion}
\label{sec:weil}

The starting point is the Weil explicit formula. For a Schwartz
function $f: \RR \to \CC$, define:
\begin{align}
\mathrm{prime\_sum}(f) &:= \sum_{p^k}
  \frac{\log p}{p^{k/2}} \bigl|\hat{f}(\log p^k)\bigr|^2, \\
\mathrm{zero\_sum}(f) &:= \sum_{\rho} \bigl|\hat{f}(\mathrm{Im}(\rho))\bigr|^2,
\end{align}
where the first sum is over prime powers $p^k$ and the second over
non-trivial zeros $\rho$ of $\zeta(s)$.

The anticommutation relation $\{D, P\} = 0$ (where $D$ is the dilation
generator and $P$ is the parity operator on $L^2(\RR)$) gives a
decomposition $L^2(\RR) = \mathcal{H}_+ \oplus \Hminus$. The
non-trivial content of the Weil criterion lies in $\Hminus$.

\begin{theorem}[Weil positivity criterion \cite{weil1952}]
\label{thm:weil}
RH is equivalent to: for all $f \in \Hminus$,
\[
W_-(f,f) := 2\bigl(\mathrm{prime\_sum}(f) - \mathrm{zero\_sum}(f)\bigr)
\geq 0.
\]
\end{theorem}

\begin{remark}
$W_-(f,f)$ is a difference of two non-negative quantities. Proving
$W_- \geq 0$ for all $f \in \Hminus$ is a universal statement over
an infinite-dimensional function space. Any proof strategy in $\MRH$
must establish this inequality.
\end{remark}

% ---- SECTION 4 ----
\section{The Model \texorpdfstring{$\MRH$}{M\_RH}}
\label{sec:model}

\begin{definition}[The model $\MRH$]
\label{def:model}
The \emph{geometric-spectral proof strategy model} $\MRH = (S, C)$
is defined as follows.

$S$ is the class of all proof strategies that proceed by one of:
\begin{itemize}
\item[\textup{(G)}] \textbf{Geometric strategy.} Construct an
  arithmetic scheme $X$ and an intersection theory $T$ on $X$,
  and derive $W_-(f,f) \geq 0$ from a Hodge-type positivity theorem
  for $T$.
\item[\textup{(S)}] \textbf{Spectral strategy.} Construct a Hilbert
  space $\mathcal{H}$ and an operator $A$ on $\mathcal{H}$, identify
  the zeta zeros $\{\gamma_n\}$ with spectral data of $A$, and derive
  $W_-(f,f) \geq 0$ from spectral properties of $A$.
\end{itemize}

$C$ consists of three constraints:
\begin{itemize}
\item[\textup{(C1)}] The strategy is formalized in ZFC.
\item[\textup{(C2)}] The strategy does not assume RH as a hypothesis.
\item[\textup{(C3)}] The strategy produces a proof of $W_-(f,f) \geq 0$
  for all $f \in \Hminus$.
\end{itemize}
\end{definition}

\begin{remark}
The model $\MRH$ is broad: it includes Arakelov geometry, $\FF_1$
geometry, tropical geometry, Weil-\'etale cohomology (geometric),
and the Hilbert--P\'olya program, Connes' scaling operator approach,
heat kernel methods, and Toeplitz operator methods (spectral). It
does not include analytic methods (zero-free regions, moment estimates)
or algebraic methods (automorphic forms, Langlands program).
\end{remark}

% ---- SECTION 5 ----
\section{The Discriminating Property \texorpdfstring{$\PAC$}{P\_AC}}
\label{sec:pac}

\begin{definition}[Arithmetic compactness property $\PAC$]
\label{def:pac}
A strategy $\mathcal{A} \in \MRH$ satisfies $\PAC$ if, at some step,
it uses a finite-dimensional \'etale cohomology group $H^1(X, \Ql)$
to encode the non-trivial zeros of $\zeta(s)$, where $X$ is a smooth
projective variety over a finite field $\FF_q$ and the
finite-dimensionality follows from the compactness (projectivity) of $X$.

Precisely: $\mathcal{A}$ satisfies $\PAC$ if it constructs $X/\FF_q$
such that the zeros of
\[
L(H^1(X), s) := \det\!\bigl(1 - q^{-s}\Frob \mid H^1(X,\Ql)\bigr)^{-1}
\]
coincide with the non-trivial zeros of $\zeta(s)$.
\end{definition}

We now verify the three conditions of Definition~\ref{def:discrim}
and Definition~\ref{def:safety}.

\subsection{Necessity}

\begin{proposition}
\label{prop:necessity}
Any strategy $\mathcal{A} \in \MRH$ that proves RH must satisfy $\PAC$,
or use a genuinely new approach outside $\MRH$.
\end{proposition}

\begin{proof}
$W_-(f,f) \geq 0$ is a universal statement over the
infinite-dimensional space $\Hminus$. To prove it, a strategy must
either:
\begin{enumerate}
\item Reduce to a finite-dimensional problem via a compact object $X$
  (satisfying $\PAC$), or
\item Establish positivity directly over the infinite-dimensional space
  without finite-dimensional reduction.
\end{enumerate}
For geometric strategies (G): the Hodge Index Theorem is a statement
about a finite-dimensional cohomology group. Applying it to $W_-$
requires identifying the infinite-dimensional space $\Hminus$ with
$H^1(X,\Ql)$ for some compact $X$. This is $\PAC$.

For spectral strategies (S): the spectral theorem for a self-adjoint
operator $A$ gives $\langle Af, f\rangle = \int \lambda\, d\langle
E_\lambda f, f\rangle$. For this to equal $W_-(f,f) =
\mathrm{prime\_sum}(f) - \mathrm{zero\_sum}(f)$, the spectral measure
must encode the zeta zeros as isolated eigenvalues of finite
multiplicity. This requires the underlying space to admit a pure point
spectrum, which requires compactness of the underlying object. This
is again $\PAC$.

Option (2)---a direct infinite-dimensional proof---is not excluded
by this argument. Such a proof would lie outside $\MRH$ by definition.
\end{proof}

\subsection{Conflict}

\begin{proposition}
\label{prop:conflict}
Any strategy $\mathcal{A} \in \MRH$ satisfying $\PAC$ cannot prove RH.
\end{proposition}

\begin{proof}
Suppose $\mathcal{A}$ satisfies $\PAC$: it constructs $X/\FF_q$ with
$L(H^1(X), s) = \zeta(s)$ (or a factor encoding all non-trivial zeros).

By the Weil conjectures (Deligne \cite{deligne1974}), for any smooth
projective $X/\FF_q$, the eigenvalues $\alpha$ of $\Frob$ on
$H^1(X, \Ql)$ satisfy $|\alpha| = q^{1/2}$. The zeros of
$L(H^1(X), s)$ are the values $s$ where $q^{-s} = \alpha^{-1}$ for
some eigenvalue $\alpha$. Since $|\alpha| = q^{1/2}$:
\[
q^{-\mathrm{Re}(s)} = |\alpha|^{-1} = q^{-1/2}
\implies \mathrm{Re}(s) = \tfrac{1}{2}.
\]
Therefore \emph{all} zeros of $L(H^1(X), s)$ lie on the critical line.

If $L(H^1(X), s) = \zeta(s)$, then all non-trivial zeros of $\zeta(s)$
lie on the critical line. That is RH.

Therefore: the identification $L(H^1(X), s) = \zeta(s)$ is equivalent
to RH. Any strategy satisfying $\PAC$ must establish this identification
as a step---but establishing it \emph{is} proving RH. The strategy
cannot prove RH without already having proved RH. Constraint (C2) is
violated.
\end{proof}

\begin{remark}
The conflict condition does not rely on any unproved conjecture. It
uses only Deligne's theorem \cite{deligne1974}, which is a theorem.
The argument is: Deligne forces $L(H^1(X),s)$ to satisfy RH; therefore
$L(H^1(X),s) = \zeta(s)$ implies $\zeta(s)$ satisfies RH; therefore
the identification step assumes what it is trying to prove.
\end{remark}

\subsection{Self-referential safety}

\begin{proposition}
\label{prop:safety}
$\PAC$ is self-referentially safe with respect to $\MRH$: no strategy
$\mathcal{A} \in \MRH$ can decide whether a given strategy satisfies
$\PAC$.
\end{proposition}

\begin{proof}
Deciding $\PAC$ for a given strategy $\mathcal{A}$ requires determining
whether $\mathcal{A}$ constructs $X/\FF_q$ with $L(H^1(X),s) = \zeta(s)$.

By Proposition~\ref{prop:conflict}, establishing $L(H^1(X),s) = \zeta(s)$
is equivalent to proving RH. Therefore, deciding $\PAC$ requires proving
RH.

But any strategy $\mathcal{A} \in \MRH$ satisfies constraint (C2): it
does not assume RH. A strategy that decides $\PAC$ would need to prove
RH as a subroutine, violating (C2).

Therefore no strategy in $\MRH$ can decide $\PAC$. \qed
\end{proof}

% ---- SECTION 6 ----
\section{The Unsatisfiability Certificate}
\label{sec:cert}

\begin{theorem}[Main theorem]
\label{thm:main}
The triple $(\MRH, f_{\mathrm{RH}}, \PAC)$ is an unsatisfiability
certificate. No strategy in $\MRH$ can prove RH.
\end{theorem}

\begin{proof}
By Propositions~\ref{prop:necessity}, \ref{prop:conflict},
and~\ref{prop:safety}, $\PAC$ is a discriminating property
(Definition~\ref{def:discrim}) that is self-referentially safe
(Definition~\ref{def:safety}). By Theorem~\ref{thm:framework},
no strategy in $\MRH$ achieves $f_{\mathrm{RH}} = 1$.
\end{proof}

\begin{corollary}
Any valid proof of RH must use a strategy outside $\MRH$: it cannot
proceed via geometric methods (Hodge Index on arithmetic schemes) or
spectral methods (operator with spectrum encoding zeta zeros) as
defined in Definition~\ref{def:model}.
\end{corollary}

\begin{remark}[What the theorem does not say]
\begin{enumerate}
\item RH is not unprovable. By Shoenfield absoluteness
  \cite{shoenfield1961}, RH is $\Pi^0_1$ and not independent of ZFC.
\item The theorem does not cover analytic or algebraic proof strategies.
\item The theorem does not identify what proof strategy \emph{would}
  work. By the Third Law of \cite{metamethod2026}, a meta-methodology
  can diagnose failure but cannot construct success.
\end{enumerate}
\end{remark}

% ---- SECTION 7 ----
\section{The Second Law: Generalization Barrier}
\label{sec:second}

The Second Law of \cite{metamethod2026} states: when a proof is
generalized from model $M_1$ to $M_2 \supset M_1$, it fails if and
only if the discriminating property $P$ becomes decidable within $M_2$.

Applied to RH:

\begin{itemize}
\item $M_1$ = geometric proof strategies for \emph{function field} RH.
  Here $\PAC$ is \emph{satisfied}: for a smooth projective curve
  $C/\FF_q$, $H^1(C, \Ql)$ is finite-dimensional (dimension $2g$),
  and Weil's proof \cite{weil1948} goes through.

\item $M_2 = \MRH$ = geometric proof strategies for \emph{number field}
  RH. Here $\PAC$ is \emph{unsatisfiable}: by
  Proposition~\ref{prop:conflict}, no compact $X/\FF_q$ has
  $L(H^1(X),s) = \zeta(s)$.
\end{itemize}

The Second Law predicts that the generalization from $M_1$ to $M_2$
fails. It does---and the failure is of a stronger type than in the
circuit complexity cases of \cite{metamethod2026}:

\begin{center}
\begin{tabular}{lll}
\toprule
Setting & Failure mode & Reason \\
\midrule
Monotone $\to$ general circuits & $P$ becomes \emph{decidable} in $M_2$
  & Model gains power to simulate $P$ \\
Function field $\to$ number field & $P$ becomes \emph{impossible} in $M_2$
  & Required object does not exist \\
\bottomrule
\end{tabular}
\end{center}

In the circuit case, the model $M_2$ can simulate the proof tool.
In the RH case, the proof tool's prerequisite (compact $X$ with
$L(H^1(X),s) = \zeta(s)$) does not exist in $M_2$. This is a
``hard'' generalization barrier: not a simulation failure but an
existence failure.

% ---- SECTION 8 ----
\section{Discussion}
\label{sec:discussion}

\subsection{Comparison with prior work}

Connes \cite{connes1999} identified the spectral interpretation of
zeta zeros and explicitly noted that the positivity of $W_-$ is the
key unproved step. The present work makes precise \emph{why} this
step cannot be proved within $\MRH$: any strategy satisfying $\PAC$
must establish $L(H^1(X),s) = \zeta(s)$, which is equivalent to RH
(Proposition~\ref{prop:conflict}).

Faltings \cite{faltings1984} proved the arithmetic Hodge Index Theorem
for single arithmetic surfaces. This does not satisfy $\PAC$ because
it does not encode the zeta zeros in $H^1$.

\subsection{The weakness in Necessity}

Proposition~\ref{prop:necessity} has a gap: it does not rule out a
direct infinite-dimensional proof of $W_- \geq 0$ that does not use
$\PAC$. Such a proof would lie outside $\MRH$ by definition, so the
certificate would not apply to it. We do not know whether such a proof
exists. This is the honest boundary of the present result.

\subsection{Open questions}

\begin{enumerate}
\item Does there exist a proof strategy outside $\MRH$ that proves RH?
  (Unknown. The certificate says nothing about this.)
\item Can the Necessity condition be strengthened to cover all proof
  strategies, not just those in $\MRH$? (This would require a much
  stronger meta-mathematical argument.)
\item Is there an analogue of $\PAC$ for analytic proof strategies?
  (Unknown. Analytic methods do not obviously require a compact object.)
\end{enumerate}

% ---- REFERENCES ----
\begin{thebibliography}{99}
\bibitem{metamethod2026}
J.~Xie, \emph{A Meta-Methodology of Unsatisfiable Proofs:
Self-Referential Safety and the Structure of Lower Bounds}, 2026.

\bibitem{deligne1974}
P.~Deligne, \emph{La conjecture de Weil. I},
Publ. Math. IHES \textbf{43} (1974), 273--307.

\bibitem{weil1948}
A.~Weil, \emph{Sur les courbes alg\'ebriques et les vari\'et\'es
qui s'en d\'eduisent}, Hermann, Paris, 1948.

\bibitem{weil1952}
A.~Weil, \emph{Sur les ``formules explicites'' de la th\'eorie des
nombres premiers}, Comm. S\'em. Math. Univ. Lund (1952), 252--265.

\bibitem{connes1999}
A.~Connes, \emph{Trace formula in noncommutative geometry and the
zeros of the Riemann zeta function},
Selecta Math. \textbf{5} (1999), 29--106.

\bibitem{faltings1984}
G.~Faltings, \emph{Calculus on arithmetic surfaces},
Ann. of Math. \textbf{119} (1984), 387--424.

\bibitem{connes-consani2014}
A.~Connes and C.~Consani, \emph{The arithmetic site},
C.~R.~Math. Acad. Sci. Paris \textbf{352} (2014), 971--975.

\bibitem{shoenfield1961}
J.~R.~Shoenfield, \emph{The problem of predicativity},
in: \emph{Essays on the Foundations of Mathematics},
Magnes Press, Jerusalem, 1961, pp.~132--139.

\bibitem{razborov-rudich1994}
A.~Razborov and S.~Rudich, \emph{Natural proofs},
J. Comput. System Sci. \textbf{55} (1997), 24--35.
\end{thebibliography}

\end{document}
